\section{Producing the warning plateaus}
\label{sec:warning-safety}
We prove the deterministic input to the warning estimate of
Section~\ref{sec:gaussian-warnings}. Fix the patch and the positive
fine-action and predictor-penalty constants $c_f,c_p$. We retain the
funded path bounds, the clean-sector precision bounds and the stationary
response of Sections~\ref{sec:sector-gram} and \ref{sec:local-mean}.
Let $t=\operatorname{fiery}(x)$, $B=\operatorname{radius}(x)$ and
$q=\dim\mathfrak g$. We use $\delta_0$ for a fixed small error allowance,
to distinguish it from the field scale $t$.

The guide estimate supplies a constant $A>0$, chosen before $L$, such
that the fine guide has norm at most $At/L^2$ and the predictor guide
at most $At$. Choose
\begin{equation}\label{eq:safety-goals}
 \delta_0=\frac{\sqrt{c_f}}{32},\qquad
 \frac{A}{L^2}\le\delta_0,\qquad
 W>\frac{2}{\sqrt{c_p}}(A+\delta_0).
\end{equation}
For example, any integer lower bound strictly larger than $A/\delta_0$
suffices for $L$, since $L\ge1$. We will choose the remaining constants
after $L$ and $W$, and finally one lower bound for $x$. None of these
choices depends on the lattice period or on the seed allocation.

\subsection{A uniform bound for the affine mean}
First we need $\|\mu_e\|\le C t$ with a fixed post-$L$ constant $C$.
A bound with upper exponent $-1/2$ alone would not imply this stronger
comparison with the particular scale $t$.

On a clean ball of fixed radius $n$, the axial gauge construction of
Section~\ref{sec:seed-geometry} bounds every coarse bond chord by
$K_n t$, where one may take
\[
                  K_n=32(2(n+1)+1)^5\#\{\text{plaquette axes}\}.
\]
The number $n$ is fixed before $x$. For a fine plaquette in the coarse
replica, use $n=4$: each of its four links has chord at most $K_4t$
in that gauge, so its curvature norm is at most $4K_4t$. Every path
term retained by the fine cutoff leaves the required radius-five margin at its endpoint inside the clean
$B$-ball: its endpoint separation plus five is at most $B$. Apply the same radius-four argument separately at that
endpoint. Unitary transport preserves its norm; summing the absolute
path coefficients bounds the entire fine base stroke by
\[
                 4K_4(\sqrt{c_f}+P_f)t,
\]
where $P_f$ is the funded absolute fine-path price. This is a sum of
individual fixed-radius estimates; no gauge on the whole growing
$B$-ball is required for this bound.

For a predictor base stroke, take $n$ equal to the fixed predictor
read radius. The logarithm chart extends the locally gauged coarse
field to a global logarithmic field of norm at most $K'_nt$, while
preserving the actual links on that stencil. Its linear replica has
no larger norm. The fixed-patch Taylor estimate from
Section~\ref{sec:predictor-branch} bounds the predictor by
$J(\|w\|+\|w\|^2)$ for the pair of logarithmic fields $w$.
Eventually $\|w\|\le K'_nt\le1$, so this is at most $2JK'_nt$.
Locality and gauge covariance identify it with the original predictor
base stroke for every original tag choice. We have therefore proved
a common bound $A_1t$ for both base-stroke types on clean balls.

Let $s>0$ be the fixed row normalization, $N_1$ the row operator
bound, and $\gamma,K$ the shell-access constants. The saved frozen
logarithmic field is bounded by $2\gamma Kt$. A selected affine row
at zero is the base stroke minus its guide plus the row applied to
this saved field. Unselected rows vanish. Hence all normalized affine
rows have norm at most $b_1t$, where
\[
                  b_1=s(A_1+A+2N_1\gamma K).
\]
Let $Y$ bound the absolute column sums of the row matrix, and let
$P_0$ bound both row and column sums of the unsymmetrized local inverse
$P$. The affine linear coefficient $\ell$ and the mean $m=-P^T\ell$
then satisfy coordinatewise
\[
                    |\ell_i|\le Yb_1t,
       \qquad |m_i|\le P_0Yb_1t,
       \qquad \|\mu_e\|\le qP_0Yb_1t.
\]
Each inequality is the elementary bound
$|\sum_jM_{ij}v_j|\le(\sum_j|M_{ij}|)\sup_j|v_j|$;
the last sums the $q$ colour coordinates. Choose once and for all
\begin{equation}\label{eq:safety-mean-cap}
                    C=1+\max\{0,b_1,Yb_1,qP_0Yb_1\}>0.
\end{equation}
The same algebra gives the residual identity
$Am+\ell=-A(P-A^{-1})^T\ell$. Its entries are bounded by a
fixed multiple of $\operatorname{distPay}(B)t$, and hence have every
prescribed negative upper exponent by the local inverse tail estimate.
The mean and its residual thus come from the same affine system.

\subsection{The genuine local minimizing branches}
We also need exact likelihood values, before restoring other masks.
Fix any $\varepsilon>0$. There exist a radius $r>0$ and a finite
read distance $d$, depending on the patch and $\varepsilon$, such that
if all fine bond chords within distance $d$ of a coarse site are at
most $r$, then its site likelihood equals one, every outgoing bond
likelihood equals one, and the predictor bonds have chord at most
$\varepsilon$.

Here is the uniformity argument. On the covering lattice, the identity
configuration lies in the open unique-minimizer branches for both
the frame and the connector. Their values there are the identity.
The continuity and uniqueness proved in
Section~\ref{sec:predictor-branch} imply a neighbourhood on which these
branches remain selected and each predictor chord is below
$\varepsilon$. Intersect these neighbourhoods over the finitely many
axes and frame-tag pairs. A product-topology neighbourhood contains
one obtained by restricting finitely many covering bonds to chord
balls of a common positive radius $r$. Enclose their anchors in a
ball of finite read distance $d$.

Pulling a periodic configuration to the cover preserves all these
finite reads, even when some periodic bonds coincide. Thus the same
$r,d$ work for every period and every tag array. On a minimizing
branch the chosen tag is uniform: its mass is the reciprocal of the
number of tags. The physical likelihood multiplies by that number,
so its value is exactly one. This proves both likelihood identities,
not merely their convergence to one.

Suppose now that the coarse field is clean on the $B$-ball about $o$
and that, on its half-ball, the actual fine field is
$U=\exp(v)\operatorname{replica}V$ with $\|v_e\|\le w(x)$,
where $w$ has upper exponent $-1/2$. The local axial gauge and logarithm
extension put both coarse and fine chords below a common bound
$W_1(x)$ of upper exponent $-1/2$. Indeed the coarse logarithm has
that bound, the exponential is locally Lipschitz, and
\[
 \operatorname{chord}(\exp(v_e)\operatorname{replica}V_e)
 \le\operatorname{chord}(\exp(v_e))
                 +\operatorname{chord}(\operatorname{replica}V_e).
\]
Choose $x$ so that $B\ge2d+12$ and
$W_1(x)\le\min\{r,\varepsilon/4\}$, applying the preceding branch
result with target $\varepsilon/2$. The predictor/coarse relative chord
is at most $\varepsilon/2+W_1(x)\le\varepsilon$. Every fine plaquette
chord is at most $4W_1(x)\le\varepsilon$. Gauge covariance transfers
these statements and the likelihood identities to the original fields.

\subsection{The three errors on a small fluctuation event}
Let $o$ have distance greater than $30B$ from the seeds. The exact
flag condition implies that the coarse field is clean on its needed
hull. For $\beta>0$, assume
\begin{equation}\label{eq:safety-watch}
 |z_i|\le\xi t\sqrt\beta
       \quad\text{for every coordinate whose anchor lies within }B
       \text{ of }o.
\end{equation}
With $v=\mu+\beta^{-1/2}z$, interpreted in the colour basis,
\eqref{eq:safety-mean-cap} gives
\[
              \|v_e-\mu_e\|\le q\xi t,
                  \qquad \|v_e\|\le(C+q\xi)t.
\]
Away from the seeds the actual fine field agrees with this increment
of the coarse replica. Thus the small-chart branch result just proved
applies to it, with $w=(C+q\xi)t$.

For a fine or predictor row based at $o$, write $B_p(U,V)$ for its
actual stroke minus its guide, and $R_p(v)$ for the transported affine
row with the same coarse field, tags and saved variables. There is a
nonnegative error $H(x)$ of upper exponent $-1$ such that
\begin{equation}\label{eq:safety-nonlinear-error}
                      \|B_p(U,V)-R_p(v)\|\le H(x).
\end{equation}
To detail the estimate, use the preceding common local gauge and
extend its two logarithmic fields by zero outside the read hull.
Let $h$ bound their norms; it has upper exponent $-1/2$. The local
product logarithm differs from the sum by at most $Jh^2$. Taylor
expansion of each genuine fine/predictor stroke at zero leaves a
remainder $E_0$ of upper exponent $-1$. Subtract the corresponding
Taylor formula at the replica. The five resulting errors are bounded
by
\begin{equation}\label{eq:safety-five-errors}
  2E_0+N_1Jh^2+N_1h e^{-\sigma R}
                           +2N_1c_GB A_2h^2.
\end{equation}
Here $R=\lfloor B/100000\rfloor$ is the retained path cutoff, and
$\sigma>0$ is the funded path decay exponent. They are, respectively,
the two Taylor remainders, the product-log
error acted on by the linear row, the cutoff tail of that row, and
the change from flat to coarse-field transport. For the last term,
each transport path has length at most $B$, each coarse chord is at
most $A_2h$, and the adjoint perturbation bound is $2c_G$ times the
path chord; multiply by the input norm $h$ and the row price $N_1$.
The retained read radius is at most $B/4$, so truncating the extended
fields does not alter the actual row. Since $B$ has upper exponent
zero and $e^{-\sigma R}$ has every negative exponent, every term in
\eqref{eq:safety-five-errors} has upper exponent $-1$. This proves
\eqref{eq:safety-nonlinear-error} with the original tags.

The affine row matrix has absolute row price $Y$ after normalization
by $s$. Therefore \eqref{eq:safety-watch} implies
\begin{equation}\label{eq:safety-linear-error}
              \|R_p(v)-R_p(\mu)\|\le q s^{-1}Y\xi t.
\end{equation}
The local mean comparison in \eqref{eq:mean-remote-final} gives a
nonnegative $E(x)$ of upper exponent $-1$ with
$\|R_p(\mu)\|\le E(x)$. Choose
\begin{equation}\label{eq:safety-xi}
 \xi=\min\left\{(q+1)^{-1},
                 \frac{\delta_0}{3q s^{-1}|Y|+1}\right\}>0.
\end{equation}
Then $q\xi\le1$ and $q s^{-1}Y\xi\le\delta_0/3$.
An upper-exponent-$-1$ function is eventually at most
$\delta_0t/3$: multiplying it by $\widehat\beta_f^{1/2}$ gives a
quantity tending to zero, while
$t=\widehat\beta_f^{-1/2}\operatorname{fineMonitor}\ge
\widehat\beta_f^{-1/2}$. Hence $H,E\le\delta_0t/3$ eventually.
The triangle inequality proves
\begin{equation}\label{eq:safety-bill}
                         \|B_p(U,V)\|\le\delta_0t.
\end{equation}

Choose the earlier small-chart target $\varepsilon$ below the fixed
curvature-comparison radius. In that neighbourhood,
$\operatorname{chord}(g)\le2\|\operatorname{curvature}(g)\|$.
Adding the guide bounds to \eqref{eq:safety-bill}, and retaining the
square-root penalty normalizations, gives
\[
 \operatorname{toll}(U,p)
       \le\frac{2}{\sqrt{c_f}}(A/L^2+\delta_0)t<t/4,
 \qquad
 \operatorname{wire}(Q(U),V,(o,a))
       \le\frac{2}{\sqrt{c_p}}(A+\delta_0)t<Wt.
\]
The strict inequalities follow from \eqref{eq:safety-goals}; the first
coefficient is at most $1/8$. Both physical likelihoods are already
exactly one by the local minimizing-branch argument.

\subsection{The cutoff margin and the final order of choices}
Set $a=\max\{C+1,\gamma K\}$. The same event has
$\|v_e\|\le(C+q\xi)t\le at$, so the broad bond cutoff is one.
For derivatives, use the event with $\xi/2$. It gives the strict
margin $C+q\xi/2<a$, while the fine and predictor inequalities above
are strict as well. Continuity makes these masks constant on a
neighbourhood of the base point. This is exactly the plateau required
in Section~\ref{sec:activity-jet-assembly}; no derivative is taken
through a chosen gauge or a varying seed assignment.

The order is now explicit: choose the guide bound $A$ and error
allowance before $L$; choose $L$, then $W$; obtain the fixed shell and
precision constants, then $C,a,\xi$; finally enlarge the $x$ threshold
for the finite-read branches, the two order-$-1$ errors, the chart
radius and the inverse-series bounds. This supplies the previously
displayed safety input on every funded, stamped sector. Funding,
the stationary response and the guide bounds retain their separate
construction hypotheses.
